
%%%%%%%%%%%%%%%%%%%%%%%%%%%%%%%%%%%
%%      Proposta de solução      %%
%%%%%%%%%%%%%%%%%%%%%%%%%%%%%%%%%%%

\section{Proposta de Solução}
    \label{sec:proposta-de-solucao}

    A seguir é elaborado o conjunto de soluções para os requisitos levantados na seção anterior e que serão implementados no sistema. As soluções levantadas são exemplos de possibilidades, baseadas em estudos de caso e em experiências anteriores. A escolha da solução final será feita durante a fase de desenvolvimento do sistema, mediante a análise de custo-benefício de cada uma das soluções.

    \subsection{Interfaces de Código}

    Uma interface de codigo pode ser descrita como um conjunto de instruções que permitem estabelecer a conexão entre duas bilbiotecas ou ate mesmo para promover o uso de uma bilbioteca. Logo o momento da definição da interface é uma das fases mais importantes da criação de um software. [REVISAR]

    Para atender os requisitos da extensabilidade, atravez da pesquisa bibliografica, as interfaces do projeto seguirão proximas as interfaces da maior bilbioteca de \textit{Machine Learning}, o scikit-learn. A biblioteca do scikit-learn é umas das mais extensas em numero de impletementaçãos de algoritmos de aprendizado de maquina, alem de possuir uma comunidade bastante ativa. isso garante que alem da compatilibidade, não so com a propria plataforma, como simplifica a extenção das funcionalidades do software.

    \subsection{Interface Gráfica}

    Para o desenvolvimento deste trabalho foi selecionado com a biblioteca baseada em Qt. O Qt provem então uma serie de impletementaçãos voltadas a interface grafica de usuario, por exeplo a criação de janelas, dialogos, entre outros. A biblioteca fornece compatilibidade com a maioria dos sistemas operacionais atuais, o que atende ao requesito não-funcional da compatilibidade como diversos sistemas operacionais.

    \subsubsection{Widgets}

	O PyQt é uma biblioteca baseada em \textit{widgets}, onde cada widget corresponde a um componente exibido na janela. Os \textit{widgets} funcionam como \textit{wrapper's}, funções que encapsulam outra chamada, diretamente para o objeto do Qt. Os \textit{widgets} do Qt são divididos em duas categorias: os \textit{widgets} de controle e os \textit{widgets} de exibição. Os \textit{widgets} de controle são aqueles que permitem a interação do usuário com o sistema, como por exemplo, botões, caixas de texto, entre outros. Já os \textit{widgets} de exibição são aqueles que exibem informações ao usuário, como por exemplo, tabelas, gráficos, entre outros.

    Os \textit{widgets} do PyQt são criados a partir de classes, são fácilmente extensíveis, logo é possível reutilizar e criar novos componentes para o sistema. A criação de novos \textit{widgets} é feita através da herança, onde uma classe é criada a partir de outra, herdando seus atributos e métodos. A herança é uma das principais características da orientação a objetos, que permite a reutilização de código, reduzindo a quantidade de código a ser escrito.

    As ações no PyQt funcionam em um \textit{event loop}, onde o sistema fica em um loop infinito, aguardando por um evento. Quando um evento é disparado, o sistema sai do loop e executa a ação correspondente ao evento. Como interfaces gráficas de usuário são orientadas a eventos, a biblioteca opera através de eventos e sinais, para isso ele utiliza o conceito de \textit{slots}. Os \textit{slots} são funções que recebem um objeto de callback para realizar uma tarefa quando um evento acontece e um sinal é enviado.

    Por exemplo, em sistema com um botão e um contador embaixo, sempre que o usuário clica no botão, o contador é incrementado. O botão envia um sinal para o slot do contador, que incrementa o contador. O sistema fica em um loop infinito, aguardando por um evento, ao clicar no botão, o sistema sai do loop e executa a ação correspondente ao evento, que é incrementar o contador.

    Logo a orientação de eventos permite que o sistema seja reponsivo para o usuário, atendendo as suas ações e apresentando uma resposta visual. No projeto, essa abordagem deve atender a praticamente todas as interações do usuário para configuração dos algoritmos de aprendizado de máquina e aos métodos de visualização.

    \subsubsection{Estilização}

    Para solucionar a questão dos estilos padrão para toda a interface e suas extensões, será utilizado a estilização por \ac{css}, que é uma linguagem de estilização que funciona a partir de uma série de regras que definem como os elementos da interface devem ser exibidos. A estilização por \ac{css} permite a criação de estilos personalizados para cada elemento da interface, permitindo a criação de estilos diferentes para cada elemento através dos seletores de elemento e de classe, como por exemplo, um botão com identificador \textit{\#BackButton}, ou a estilização de caixa de texto com a classe \textit{.Entrada}, entre outros.

    Outra possibilidade é a utilização de estilos pré-definidos, como por exemplo, o \textit{QDarkStyle} já presente na biblioteca, que fornece um estilo escuro para a interface, podendo ser utilizado para melhorar a experiência do usuário, principalmente em ambientes com pouca iluminação. Além disso também é possível utilizar biliotecas próprias para estilização da interface, que já dispõe de estilos conhecidos como o \textit{Material Design}, que é um estilo de interface que utiliza elementos como por exemplo, sombras, bordas arredondadas, entre outros.

    A estilização da interface é uma etapa importante para criar uma interface agradável para o usuário e deve ser feita de forma que corresponda às necessidades do usuário, como usar cores facilmente reconhecíveis com contraste, saturação e brilho suficientes, além de usar elementos reconhecíveis , como botões com bordas arredondadas, etc. Tudo isso oferece aos usuários uma interface intuitiva e agradável.

    \subsubsection{Gráficos}

    Gráficos são uma das principais formas de visualizar dados e são uma das formas mais comuns de apresentar informações aos usuários. A biblioteca \textit{Matplotlib} e sua extensão \textit{Seaborn} são usadas para criar gráficos que são exibidos na interface via \textit{widget} \textit{MplWidget}. Este \textit{weidget} é uma classe herdada de \textit{FigureCanvas} e \textit{NavigationToolbar}, que são classes da biblioteca \textit{Matplotlib} responsáveis por criar o gráfico e a barra de ferramentas, respectivamente.

    \textit{Matplotlib} é uma biblioteca gráfica que permite a criação de gráficos 2D e 3D, além de gráficos animados. Esta biblioteca é usada para criar gráficos de linhas, gráficos de barras, histogramas, etc. Já a biblioteca \textit{Seaborn} é uma extensão do \textit{Matplotlib} que permite a criação de gráficos melhores, com melhores estilos e cores, e além disso permite a criação de gráficos mais complexos, como gráficos de dispersão, intervalo de gráficos de linha confiável e assim sobre. A biblioteca \textit{Seaborn} deve atender as necessidades de visualização de dados do projeto, pois permite a criação de gráficos mais complexos e com melhor estilização, além de que poderá ser facilmente extensível para a criação de novos gráficos.

    \subsection{Gerenciamento de Dados}

    Para solucionar a questão de gerenciamento de dados, será utilizado o \textit{Pandas} para a manipulação dos dados, que é uma biblioteca de código aberto que fornece estruturas de dados e ferramentas de análise de dados para a linguagem de programação Python. O \textit{Pandas} é uma biblioteca que permite a manipulação de dados, como por exemplo, a leitura de arquivos, a criação de tabelas, a criação de gráficos, entre outros. A biblioteca \textit{Pandas} é uma das bibliotecas mais utilizadas para manipulação de dados, pois permite a manipulação de dados de forma simples e eficiente, além de que possui uma grande comunidade de usuários e desenvolvedores.

    O \textit{Pandas} utiliza o núcleo do \textit{Numpy}, que é uma biblioteca de código aberto que fornece suporte para arrays e matrizes multidimensionais, além de funções matemáticas para operações com esses arrays. A perfomance do \textit{Numpy} é devido ao uso da linguagem C para a criação e operação com arrays, garantiu uma grande velocidade de processamento, além de que o \textit{Numpy} é uma biblioteca que é amplamente utilizada em outras bibliotecas, como por exemplo, o \textit{Tensorflow}.

    Utilizando o \textit{Pandas} é possível importar e exportar tabelas facilmente, já que a bilbiotecas já possui esses métodos. Além disso também é possível realizar operações como ordenação, filtragem, etc, tarefas básicas da análise de dados.
